\section{模型评价与推广}
\subsection{模型优点}
模型显式分离瞬时容量状态与小时重叠电量，避免 fractional-duration 任务被按整小时结算。Q1 的预测、基础调度及图形来源保持独立；Q2 五个 direct 方向共享语义唯一的 BaseCandidate 池与 6 遍相同顺序，直接碳与多起点强化分离，平衡评分使用公共 ideal--nadir 尺度；Q3/Q4 通过哈希绑定 formal balanced 下游输入，并共享来源显式能源流、净电网指标和可用 SOC 区间 EFC。增量/完整重算、全量硬约束、负向篡改、双运行和包清单共同构成可复核证据链。

五个 direct 方向均覆盖 50000 个任务，七个 Q3 策略和 17 个 Q4 情景均有逐行结构化结果。双运行确定性 bundle 哈希一致，故当前结果达到无独立 Oracle 条件下的 scenario-feasible 可信等级。

\subsection{模型不足}
Q2 已在共同第 5 遍达到预声明的主目标/接受率阈值稳定，但第 6 遍成本和平衡仍分别接受 4 和 1 个极小改善，且每任务最多只有三个基础候选；因此它不是零接受固定点，更没有全局最优证书。Q4 每轮只检查 600 个候选并最多 8 轮，同样只是有限预算 best-found。Q3 的 21 状态 SOC 也只给出网格内候选。储能未计退化、自放电和备用容量；所选方案 EFC 为 92.401，仍不能把账面收益直接解释为可部署收益。

通信侧只使用静态单向时延，未纳入带宽、拥塞和迁移能耗。formal balanced 的服务目标为 131.4759、服务质量为 0.007549，最大等待达到 1926 小时，说明综合评分没有消除尾部服务风险。Q4 正常联合基准购电碳为 0；在新能源短缺压力参照下，0\% 和 10\% 档找到达标候选，20\% 和 30\% 档只是在声明候选银行内未找到达标轨迹，并非数学不可行。负成本同样主要来自题面售电边界。

\subsection{问题一的鲁棒性与可追溯性}
题面单窗口切分被保留，并以 8 个完全位于最终测试之前的 expanding-window 原点检查时间稳定性。最终测试三种 WAPE 可由 432 条预测逐项复算，8 个滚动窗口的三口径曲线可由 3456 条选定候选逐小时预测复算；\texttt{schedule\_provenance.json} 绑定 Q1 全量表、538 条末端表、Q2 balanced 表及两张 Q1 图。负向替换测试会在 Q1/Q2 哈希同源或 provenance 漂移时失败。

System Aggregate WAPE 明显低于 Macro 与 Micro，部分来自序列间误差抵消，不能解释单序列容量风险。Q1 调度虽通过全时域审计，仍是确定性贪心的 scenario-feasible incumbent。

\subsection{推广方向}
后续应扩大 Q2 每任务时空邻域、比较零接受停止与阈值停止，并引入混合整数规划上下界，并继续使用同一公共候选与全量审计。Q3 可加入每 MWh 吞吐退化成本和年度循环上限；在线部署可采用滚动时域并传递绝对 SOC。多目标部分可用 $\varepsilon$-约束法扩充非支配候选\cite{miettinen1999nonlinear}，但仍须区分“候选内非支配”和“已证 Pareto”。
